% 长方体三视图：左侧画立体，右侧按位置摆放主视/俯视/左视
\documentclass[tikz,border=4pt]{standalone}
\usepackage{ctex}
\usepackage{amsmath}
\usetikzlibrary{calc, arrows.meta}
\begin{document}
\begin{tikzpicture}[thick, scale=1.0,
    vlabel/.style={font=\footnotesize}]
  % --- 左：长方体 a=3, b=1.5, c=2 ---
  % 透视：用 (a, 0) 作为长方向、(0,c) 为高、({b*cos(30)}, {b*sin(30)}) 为宽
  \pgfmathsetmacro{\a}{3}
  \pgfmathsetmacro{\b}{1.5}
  \pgfmathsetmacro{\c}{2}
  \pgfmathsetmacro{\dx}{\b*cos(30)}
  \pgfmathsetmacro{\dy}{\b*sin(30)}
  \coordinate (A) at (0,0);
  \coordinate (B) at (\a,0);
  \coordinate (C) at (\a+\dx, \dy);
  \coordinate (D) at (\dx, \dy);
  \coordinate (A1) at (0,\c);
  \coordinate (B1) at (\a,\c);
  \coordinate (C1) at (\a+\dx, \c+\dy);
  \coordinate (D1) at (\dx, \c+\dy);
  % 看得见的边
  \draw (A) -- (B) -- (B1) -- (A1) -- cycle;
  \draw (B) -- (C) -- (C1) -- (B1);
  \draw (A1) -- (D1) -- (C1);
  % 看不见的边（虚线）
  \draw[dashed, gray] (A) -- (D) -- (C);
  \draw[dashed, gray] (D) -- (D1);
  % 方向箭头：主视（从正前方）、俯视（从上方）、左视（从左侧）
  \draw[->, blue, thick] (\a/2, -1.2) -- (\a/2, -0.2) node[midway, right, font=\footnotesize, blue] {主视};
  \draw[->, blue, thick] (-1.2, \c/2 + \dy/2) -- (-0.2, \c/2 + \dy/2) node[midway, above, font=\footnotesize, blue] {左视};
  \draw[->, blue, thick] (\a/2 + \dx/2, \c+\dy+1.2) -- (\a/2 + \dx/2, \c+\dy+0.2) node[midway, right, font=\footnotesize, blue] {俯视};
  \node[below, font=\footnotesize] at (\a/2, -1.6) {（a）立体};

  % --- 右：三视图按标准位置摆放 ---
  \begin{scope}[xshift=6cm, yshift=0cm]
    % 主视（左上）：长 a × 高 c
    \draw (0, 0) rectangle (\a, \c);
    \node[font=\footnotesize] at (\a/2, \c/2) {主视};
    % 左视（右上）：宽 b × 高 c
    \draw (\a + 0.6, 0) rectangle (\a + 0.6 + \b, \c);
    \node[font=\footnotesize] at (\a + 0.6 + \b/2, \c/2) {左视};
    % 俯视（左下）：长 a × 宽 b
    \draw (0, -0.6 - \b) rectangle (\a, -0.6);
    \node[font=\footnotesize] at (\a/2, -0.6 - \b/2) {俯视};
    % 对应虚线（长对正、高平齐）
    \draw[dashed, gray, very thin] (0,0) -- (0,-0.6);
    \draw[dashed, gray, very thin] (\a,0) -- (\a,-0.6);
    \draw[dashed, gray, very thin] (\a,0) -- (\a + 0.6, 0);
    \draw[dashed, gray, very thin] (\a,\c) -- (\a + 0.6, \c);
    \node[below, font=\footnotesize] at (\a/2 + 0.3, -0.6 - \b - 0.6) {（b）三视图};
  \end{scope}
\end{tikzpicture}
\end{document}
